\chapter{Phụ lục giao diện người dùng}
\label{app:ui}

Giao diện web cho phép chơi người--máy, máy--máy và xem lại ván từ log. Phụ lục này mô tả cách
chạy, các endpoint và cách đọc dữ liệu hiển thị. Mục đích là để người đọc có thể tự kiểm tra
hành vi của tác tử, không chỉ đọc bảng.

\section{Chạy giao diện}
\begin{lstlisting}[numbers=none]
uv sync --group dev --extra ui --extra nn
uv run uvicorn caro.ui.server:app --reload
\end{lstlisting}
Mở \url{http://127.0.0.1:8000}. Khi sửa code, \texttt{--reload} tự khởi động lại. Đổi cổng bằng
\texttt{--port 8001}. Cấp neural cần đúng checkpoint; nếu model chưa tải, giao diện báo lỗi
thay vì treo.

\section{Chức năng}
\begin{itemize}
  \item \textbf{Người--máy:} chọn cấp cho một hoặc hai bên; bấm ô để đi; nút ``AI đi'' gọi tác
        tử. Nước bất hợp lệ bị từ chối và bàn cờ không đổi.
  \item \textbf{Máy--máy:} chọn hai tác tử, bấm chạy; mỗi nước hiển thị thời gian, số nút hoặc
        số mô phỏng, biến chính và bản đồ nhiệt lượt thăm (nếu có).
  \item \textbf{Xem lại:} đọc \texttt{runs/**/games.jsonl}, dựng lại ván theo đúng luật; có thể
        tua từng nước.
  \item \textbf{Lùi nước:} lùi một hoặc nhiều nước, không lùi qua mốc khai cuộc.
  \item \textbf{Ngân sách:} chọn thời gian, độ sâu hoặc số mô phỏng trong giới hạn an toàn.
\end{itemize}

\section{Endpoint chính}
\begin{center}\small
\begin{tabularx}{\textwidth}{@{}l l X@{}}
\toprule
\textbf{Phương thức} & \textbf{Đường dẫn} & \textbf{Ý nghĩa} \\
\midrule
GET & \texttt{/agents} & Danh sách tác tử, cấp, kind, khả dụng, ngân sách \\
GET & \texttt{/budgets} & Ngân sách mặc định và giới hạn \\
GET & \texttt{/openings} & Danh sách khai cuộc \\
POST & \texttt{/games} & Tạo ván mới (người--máy hoặc máy--máy) \\
GET & \texttt{/games/\{id\}} & Trạng thái ván, lịch sử, thống kê từng nước \\
POST & \texttt{/games/\{id\}/move} & Người đi một nước \\
POST & \texttt{/games/\{id\}/ai-move} & Gọi tác tử đi một nước \\
POST & \texttt{/games/\{id\}/undo} & Lùi nước \\
GET & \texttt{/replays} & Liệt kê ván trong log \\
GET & \texttt{/replays/\{id\}} & Dựng lại một ván \\
\bottomrule
\end{tabularx}
\end{center}

\section{Cách đọc thống kê}
Khi tác tử đi, giao diện hiển thị \texttt{last\_stats()}. Các khoá có thể có:
\begin{itemize}
  \item \texttt{nodes}, \texttt{depth}, \texttt{pv}, \texttt{score}: Alpha-Beta.
  \item \texttt{sims}, \texttt{visits}, \texttt{pv}: MCTS/PUCT. \texttt{visits} là bản đồ nhiệt
        số lượt thăm từng ô.
  \item \texttt{policy}, \texttt{value}: policy-only/PUCT. \texttt{policy} là xác suất các ô;
        \texttt{value} là giá trị theo góc nhìn bên đi.
  \item \texttt{reason}: cấp 2, cho biết luật đã dùng (P1--P7).
\end{itemize}
Cần phân biệt \emph{giải thích quá trình} với \emph{chứng minh tối ưu}. Một ô có nhiều lượt
thăm hoặc prior cao không chứng minh nước đó tốt; nó chỉ cho biết tác tử đã dành ngân sách cho
ô đó. Muốn đánh giá chất lượng, phải đấu và phân tích ván.

\section{Giới hạn}
Giao diện giữ ván trong bộ nhớ tiến trình; tắt server là mất ván đang chơi. Ván đã ghi trong
\texttt{games.jsonl} vẫn xem lại được. Số ván tối đa trong bộ nhớ là 256; ván cũ nhất bị bỏ.
Một nước đang chạy trong thread không huỷ được, nên ngân sách bị chặn trước (thời gian tối đa
60 giây, độ sâu 8, mô phỏng 100\,000). Các giới hạn này để tránh treo server.

\section{Kiểm thử giao diện}
Bộ test \texttt{tests/test\_ui.py} kiểm tra tạo ván, đi nước hợp lệ/bất hợp lệ, lùi nước, gọi
AI, xem lại và các trường hợp lỗi. Test dùng FastAPI TestClient, không cần chạy server thật.